\section{RL as probabilistic inference：有用的解释，不是万能证明}

讲者随后把 Decision Transformer 放进“control as inference”视角：把高 reward 解释为轨迹具有较高 optimality probability，再通过条件推断选择动作。这个视角帮助理解 return conditioning，却不能替代离线数据覆盖、因果识别或真实环境评价。

\subsection{Optimality variable 的直觉}

引入二元 optimality variable $\mathcal{O}_t$，用 reward 定义
\[
p(\mathcal{O}_t=1\mid s_t,a_t)\propto \exp\left(\frac{r(s_t,a_t)}{\alpha}\right),
\]
其中 $\alpha>0$ 是 temperature。高 reward 的 state-action 对更可能被视为 optimal。整条轨迹在 ``optimal'' 条件下的后验可写成 prior dynamics 与 reward weighting 的组合；policy learning 可理解为近似该后验。

\lecturefigure{slide-13-rl-as-probabilistic-inference.jpg}{RL as probabilistic inference：用 optimality variables 将 reward 转成条件推断问题。}{Stanford Online 官方视频 30:00；右图引用 Levine 2018。}

\begin{knowledgebox}{读图：Decision Transformer 与 inference view 的连接}
左侧仍是 $(R,s,a)$ token sequence；右侧 probabilistic graphical model 把每一步的 optimality variable 与 state/action 相连。课程的连接是：return-to-go 可以被看作对未来 optimality 的压缩条件，Transformer 则学习给定该条件的 action distribution。它不是精确执行图模型推断，而是用监督数据近似相关条件分布。
\end{knowledgebox}

\subsection{这个视角能解释什么}

它解释了为什么“条件化在高回报”可能偏向高质量动作，也解释了为什么 behavior prior 很重要：如果某种动作从未出现在日志里，posterior approximation 很难可靠支持它。它还提示 temperature、reward scale 与 entropy regularization 之间的关系，帮助连接 maximum-entropy RL。

\subsection{这个视角不能证明什么}

概率解释不自动证明目标 return 可达，也不证明条件模型能在 dataset support 外正确外推；更不证明 Transformer 已经完成因果 planning。模型学习的是 observational trajectory distribution 中的相关结构，环境动力学与 action consequence 仍需通过数据体现。

\begin{warningbox}{相关轨迹模式不是反事实保证}
若高 return 轨迹总伴随某个无因果作用的视觉特征，模型可能把它当作行为线索。只有在 intervention、环境随机性、跨 policy 数据或专门 counterfactual evaluation 下，才能更有把握地区分真正的 action consequence 与日志相关性。
\end{warningbox}

\subsection{本章小结}

RL-as-inference 为 return conditioning 提供了统一直觉：高 reward 轨迹获得更高 optimality weight，policy 近似条件后验。但该解释仍受 behavior prior 和数据支持约束，不能替代闭环实证。

